% Readings, Section B: Extracts from elementary algebra and analysis
% Page references in vocabulary notes are to the 1972 Springer edition.
\newcommand{\bcorrnote}[1]{\footnote{{\fontspec{Noto Serif}\textbf{校勘：}#1}}}

\clearpage
\begin{center}
{\sffamily\color{BellaBlue}\Large SECTION B\par}
\vspace{0.8ex}
{\bfseries\Large EXTRACTS FROM ELEMENTARY ALGEBRA AND ANALYSIS\par}
\end{center}
\addcontentsline{toc}{section}{Section B: Elementary algebra and analysis}
\markright{SECTION B\quad Elementary algebra and analysis}
\begingroup
\sloppy
\setlength{\emergencystretch}{4em}

\readhead{B1}{Операции над множествами}{Operations on sets}
\begin{rusblock}
Объединением $A\vee B$ множеств $A$ и $B$ называется множество элементов, принадлежащих либо $A$, либо $B$, либо $A$ и $B$. Пересечением $A\wedge B$ называется множество элементов, принадлежащих и $A$ и $B$. Разностью $A\backslash B$ называется множество элементов $A$ и не $B$.\bcorrnote{底本此处作“Разность $A\backslash B$ называется множество...”，格支配有误；按俄语句法改为“Разностью ... называется множество...”。}
\end{rusblock}
\vocabline{\textbf{объединение} 89 \emph{union} (один- uni- one).}

\readhead{B2}{Свойства операций над множествами}{Properties of the operations on sets}
\begin{rusblock}
\begin{enumerate}[leftmargin=2.2em]
\item Коммутативность объединения:
\[A\vee B=B\vee A.\]
\item Ассоциативность объединения:
\[A\vee(B\vee C)=(A\vee B)\vee C.\]
\item Коммутативность пересечения:
\[A\wedge B=B\wedge A.\]
\item Ассоциативность пересечения:
\[A\wedge(B\wedge C)=(A\wedge B)\wedge C.\]
\item Дистрибутивность пересечения относительно объединения:
\[A\wedge(B\vee C)=(A\wedge B)\vee(A\wedge C).\]
\item Дистрибутивность объединения относительно пересечения:
\[A\vee(B\wedge C)=(A\vee B)\wedge(A\vee C).\]
\end{enumerate}
\end{rusblock}
\vocabline{\textbf{свойство} 45 \emph{property}.}

\readhead{B3}{Взаимно однозначное соответствие}{One-to-one correspondence}
\begin{rusblock}
Взаимно однозначным соответствием между множествами $X$ и $Y$ называется соответствие, имеющее следующие свойства:
\begin{enumerate}[leftmargin=2.2em]
\item каждому элементу множества $X$ соответствует один и только один элемент множества $Y$;
\item различным элементам множества $X$ соответствуют различные элементы множества $Y$;
\item всякому элементу множества $Y$ соответствует по меньшей мере один элемент множества $X$.
\end{enumerate}
Первые два свойства определяют взаимно однозначное отображение множества $X$ на некоторое подмножество множества $Y$. В этом случае говорят о взаимно однозначном отображении $X$ в $Y$.
\end{rusblock}
\bcorrnote{底本第3项作“соответствует по меньшей мере одному элементу”，其中作为“соответствует”主语的数量短语应为主格；故改为“по меньшей мере один элемент”。}
\vocabline{\textbf{однозначный} 89, 96 \emph{one-valued} (один- one; зна- know); \textbf{мера} 99 \emph{measure} (мер-); \textbf{по меньшей мере} \emph{at least} (lit. by least measure); \textbf{отображение} 101 \emph{mapping} (рез- cut).}

\readhead{B4}{Эквивалентные множества}{Equivalent sets}
\begin{rusblock}
\textbf{Определение.} Два множества $X$ и $Y$, между которыми можно установить взаимно однозначное соответствие, называются эквивалентными, что обозначается символом $X\sim Y$. Об эквивалентных множествах говорят, что они имеют одинаковую мощность.

Соотношение эквивалентности обладает следующими свойствами:
\begin{enumerate}[leftmargin=2.2em]
\item рефлексивностью: $X\sim X$;
\item симметрией: если $X\sim Y$, то и $Y\sim X$;
\item транзитивностью: если $X\sim Y$ и $Y\sim Z$, то $X\sim Z$.
\end{enumerate}
\end{rusblock}
\bcorrnote{底本反身性一项误排为 $X\sim Y$；反身性应为 $X\sim X$，此处直接改正。}
\vocabline{\textbf{мощность} 100 \emph{power, cardinality} (мочь- be able).}

\readhead{B5}{Упорядоченные множества}{Ordered sets}
\begin{rusblock}
\textbf{Определение.} Множество $M$ называется упорядоченным, если между его элементами установлено некоторое отношение $a<b$, обладающее следующими свойствами:
\begin{enumerate}[leftmargin=2.2em]
\item между любыми элементами $a$ и $b$ существует одно и только одно из соотношений $a=b$, $a<b$, $b<a$;
\item для любых элементов $a,b,c$ из $a<b$, $b<c$ следует $a<c$.
\end{enumerate}
\end{rusblock}
\vocabline{\textbf{упорядочить} 106 \emph{to order} (ряд- row); \textbf{обладать} 105 \emph{possess} (влад-).}

\readhead{B6}{Подобные множества}{Similar sets}
\begin{rusblock}
\textbf{Определение.} Упорядоченные множества $A$ и $B$ называются подобными, если между ними можно установить взаимно однозначное соответствие, такое, что из $a_1\to b_1$, $a_2\to b_2$ и $a_1<a_2$ следует $b_1<b_2$. О подобных множествах говорят, что они имеют один и тот же тип. Отношение подобия множеств $A$ и $B$ обозначается так: $A\approx B$.

Подобные множества эквивалентны: из $A\approx B$ следует $A\sim B$. Если конечные упорядоченные множества эквивалентны, то они подобны.
\end{rusblock}
\vocabline{\textbf{подобный} 105 \emph{similar}; \textbf{подобие} 105 \emph{similarity} (доб- suitable).}

\readhead{B7}{Алгебраические операции}{Algebraic operations}
\begin{rusblock}
\textbf{Определение.} Такое соответствие, что каждой упорядоченной паре $a,b$ элементов множества $M$ соответствует единственный элемент $c$ того же множества $M$, называется алгебраической операцией, определённой в $M$.

Другими словами, алгебраическая операция, определённая в множестве $M$, есть функция, определённая на множестве всех упорядоченных пар элементов $M$, значения которой принадлежат $M$.
\end{rusblock}
\vocabline{\textbf{единственный} 89 \emph{unique} (один- uni- one); \textbf{слово} 107 \emph{word}.}

\readhead{B8}{Кольца}{Rings}
\begin{rusblock}
Непустое множество $R$ называется кольцом, если в нём определены две алгебраические операции: сложение, ставящее в соответствие каждым двум элементам $a$ и $b$ элемент $a+b$, называемый их суммой, и умножение, ставящее в соответствие каждым двум элементам $a$ и $b$ элемент $ab$, называемый их произведением, причём эти операции обладают следующими свойствами:
\begin{enumerate}[leftmargin=2.2em]
\item коммутативностью сложения: $a+b=b+a$;
\item ассоциативностью сложения: $a+(b+c)=(a+b)+c$;
\item обратимостью сложения: для любых $a$ и $b$ из $R$ уравнение $a+x=b$ имеет решение, т.е. существует элемент $c\in R$ такой, что $a+c=b$;
\item коммутативностью умножения: $ab=ba$;
\item ассоциативностью умножения: $a(bc)=(ab)c$;
\item дистрибутивностью умножения относительно сложения: $(a+b)c=ac+bc$.
\end{enumerate}
\end{rusblock}
\vocabline{\textbf{кольцо} 107 \emph{ring}; \textbf{непустой} 49 \emph{non-empty}; \textbf{сложение} 82 \emph{addition} (лаг- lay); \textbf{умножение} 88 \emph{multiplication} (мног- mult- many).}

\readhead{B9}{Примеры колец}{Examples of rings}
\begin{rusblock}
\begin{enumerate}[leftmargin=2.2em]
\item Множество целых, рациональных, действительных, комплексных чисел.
\item Множество, состоящее из одного числа $0$.
\item Множество комплексных чисел $a+bi$ с целыми $a$ и $b$ --- так называемое кольцо целых гауссовских чисел.
\item Множество действительных чисел вида $a+b\sqrt2$, где $a$ и $b$ --- целые числа.
\end{enumerate}
\end{rusblock}
\vocabline{\textbf{состоять} 85 \emph{consist of} (ста- sist- set up); \textbf{так называемый} \emph{so-called}.}

\readhead{B10}{Нуль кольца}{The zero of a ring}
\begin{rusblock}
\textbf{Определение.} Нулём кольца $R$ называется такой элемент $0$, который удовлетворяет условиям
\[
a+0=0+a=a
\]
для любого элемента $a$ из $R$.

\textbf{Теорема.} Во всяком кольце существует нуль и притом только один.
\end{rusblock}
\vocabline{\textbf{притом} 50 \emph{moreover}.}

\readhead{B11}{Области целостности}{Domains of integrity}
\begin{rusblock}
\textbf{Теорема.} Если один из сомножителей равен нулю, то и произведение равно нулю, т.е. для любого элемента $a$ кольца $R$ имеем
\[
a\cdot0=0\cdot a=0.
\]
Однако обратная теорема не верна для любых колец. Например, пары $(a,b)$ целых чисел образуют кольцо, если операции определены по формулам
\[
(a,b)+(c,d)=(a+c,b+d),\qquad (a,b)(c,d)=(ac,bd),
\]
и в этом кольце является нулём пара $(0,0)$. Если $a\ne0$ и $b\ne0$, то пары $(a,0)$ и $(0,b)$ отличны от нуля кольца, но
\[
(a,0)(0,b)=(0,0).
\]
Элементы $a$ и $b$ кольца, для которых $a\ne0$, $b\ne0$, но $ab=0$, называются делителями нуля. Кольцо без делителей нуля называется областью целостности.
\end{rusblock}
\vocabline{\textbf{целостность} 107 \emph{wholeness, integrity}; \textbf{однако} 50 \emph{nevertheless}; \textbf{делитель} 95 \emph{divisor} (дел-); \textbf{без} 52 \emph{without}.}

\readhead{B12}{Поля}{Fields}
\begin{rusblock}
Полем называется кольцо $P$, обладающее следующими свойствами:
\begin{enumerate}[leftmargin=2.2em]
\item обратимостью умножения: для любых $a$ и $b$ из $P$, где $a\ne0$, уравнение $ax=b$ имеет решение, т.е. существует элемент $q\in P$ такой, что $aq=b$;
\item $P$ содержит по меньшей мере один элемент, отличный от нуля.
\end{enumerate}
Примеры полей:
\begin{enumerate}[leftmargin=2.2em]
\item множество рациональных, действительных, комплексных чисел;
\item множество комплексных чисел $a+bi$ с любыми рациональными $a,b$ (так называемое поле гауссовых чисел);
\item множество действительных чисел вида $a+b\sqrt2$ с любыми рациональными $a$ и $b$;
\item множество из двух элементов, которые мы обозначим через $0$ и $1$, при следующем определении операций:
\[
0+0=1+1=0,\qquad 0+1=1+0=1,
\]
\[
0\cdot0=0\cdot1=1\cdot0=0,\qquad 1\cdot1=1.
\]
\end{enumerate}
\end{rusblock}
\vocabline{\textbf{поле} 107 \emph{field}; \textbf{обратимость} 92 \emph{invertibility} (верт- vert- turn).}

\readhead{B13}{Единица}{Unit element}
\begin{rusblock}
\textbf{Определение.} Единицей кольца (в частности, поля) называется такой элемент $1$, который удовлетворяет условиям
\[
a\cdot1=1\cdot a=a
\]
для любого элемента $a$.

\textbf{Теорема.} Во всяком поле существует единица и притом только одна. При этом $1\ne0$. Если кольцо содержит по меньшей мере одну единицу, то она будет единственной единицей; однако существуют кольца без единицы, например кольцо целых чисел вида
\[
\ldots,-2n,-n,0,n,2n,\ldots,\qquad n>1.
\]
\end{rusblock}
\vocabline{\textbf{частность} 107 \emph{detail, particularity} (част- part); \textbf{в частности} \emph{in particular}.}

\readhead{B14}{Деление}{Division}
\begin{rusblock}
\textbf{Теорема.} Умножение в поле обладает единственной обратной операцией, определённой для любой пары элементов, в которой второй элемент отличен от $0$; другими словами, уравнения $ax=b$ и $ya=b$ при $a\ne0$ имеют единственное решение. Обратная операция для умножения называется делением и обозначается через $b/a$.
\end{rusblock}

\readhead{B15}{Характеристика поля; простые поля}{Characteristic of a field; prime fields}
\begin{rusblock}
\textbf{Определение.} Характеристикой поля $P$ называется число $0$, если
\[
na\equiv a+a+\cdots+a\ne0
\]
для любого элемента $a\ne0$ и любого целого числа $n\ne0$; но если имеется (простое) число $p$, такое, что $pa=0$ для любого элемента $a$, то характеристикой поля $P$ называется число $p$.

Так как для числа $1$ и любого целого $n$ будет $n\cdot1=n$, то все числовые поля имеют характеристику $0$.

Примером поля характеристики $p$ (для любого простого $p$) является кольцо вычетов по модулю $p$.

\textbf{Определение.} Поле, не имеющее собственных подпольев, называется простым.

\textbf{Теорема.} Любое поле содержит простое подполе и притом только одно.
\end{rusblock}
\vocabline{\textbf{простой} 106 \emph{simple, prime} (прост-); \textbf{числовой} 103 \emph{numerical} (чит- reckon); \textbf{вычет} 104 \emph{residue} (чит- reckon); \textbf{собственный} 45 \emph{proper}; \textbf{подполе} 107 \emph{subfield}.}

\readhead{B16}{Изоморфизм}{Isomorphism}
\begin{rusblock}
\textbf{Определение.} Два множества $M$ и $M'$, в каждом из которых определена (через некоторые условия) система отношений $S$, называются изоморфными (в символах $M\cong M'$) относительно данной системы отношений, если между ними существует взаимно однозначное соответствие, такое, что если любые элементы множества $M$ находятся в любом из отношений системы $S$, то соответствующие элементы множества $M'$ находятся в том же отношении, и обратно.

Например, в случае отсутствия отношений определение изоморфизма обращается в определение эквивалентности, а в случае одного отношения $a<b$ --- в отношение подобия.
\end{rusblock}
\vocabline{\textbf{отсутствие} 72 \emph{absence}.}

\readhead{B17}{Расположенные кольца}{Ordered rings}
\begin{rusblock}
С отношением порядка в кольце связаны понятия положительности, отрицательности и абсолютной величины элементов.

\textbf{Определение.} Кольцо (в частности, поле) $R$ называется расположенным, если для его элементов определено свойство быть положительным, удовлетворяющее следующим условиям:
\begin{enumerate}[leftmargin=2.2em]
\item для любого элемента $a\in R$ имеет место одно и только одно из соотношений: $a=0$, $a$ положителен, $-a$ положителен;
\item если $a$ и $b$ положительны, то $a+b$ и $ab$ также положительны.
\end{enumerate}
\end{rusblock}
\vocabline{\textbf{расположить} 82 \emph{dispose, arrange, order} (лаг- pos- lay); \textbf{связать} 105 \emph{to connect} (вяз- knot); \textbf{положительность} 82 \emph{positivity} (лаг- pos- lay); \textbf{отрицательность} 106 \emph{negativity} (реч- speak); \textbf{место} 106 \emph{place, locus} (мест-).}

\readhead{B18}{Свойства расположенных колец}{Properties of ordered rings}\bcorrnote{底本英文小标题将 \emph{Properties} 误排作 \emph{porperties}；此处改正。}
\begin{rusblock}
\textbf{Теорема.} Если в расположенном кольце $R$ определять порядок, считая $a>b$ тогда и только тогда, когда $a-b$ положителен, то $R$ будет упорядоченным множеством, причём нуль будет меньше всех положительных и больше всех отрицательных элементов.

\textbf{Теорема.} Расположенное кольцо не имеет делителей нуля.

\textbf{Теорема.} Характеристика расположенного поля $P$ равна нулю.

\textbf{Теорема.} Сумма квадратов (и, в частности, всякий квадрат) конечного числа элементов расположенного кольца больше или равна нулю, причём равенство имеет место только в том случае, когда все данные элементы равны нулю.

\textbf{Определение.} Расположенное кольцо называется плотным, если для любых двух различных элементов $a$ и $b$ существует элемент, лежащий между $a$ и $b$, т.е. такой, что или $a<c<b$, или $b<c<a$.

\textbf{Теорема.} Всякое расположенное поле плотно.
\end{rusblock}
\vocabline{\textbf{больше} 49 \emph{greater}; \textbf{плотный} 106 \emph{dense} (плот- thick).}

\readhead{B19}{Аксиома Архимеда}{Axiom of Archimedes}
\begin{rusblock}
Есть свойство чисел, которое не переносится на любые расположенные кольца. Это --- выполнение так называемой аксиомы Архимеда.

\textbf{Определение.} Кольцо (в частности, поле) называется архимедовски расположенным, если оно обладает следующим свойством:

\textbf{Аксиома Архимеда:} для любых элементов $a$ и $b$ кольца, где $b>0$, существует натуральное число $n$ такое, что $nb>a$.

\textbf{Пример неархимедовски расположенного кольца.} Пусть $R$ есть кольцо многочленов
\[
f(x)=a_0+a_1x+a_2x^2+\cdots+a_nx^n
\]
с рациональными коэффициентами (при обычных операциях сложения и умножения). Будем считать многочлен положительным, если его старший коэффициент $a_n$ положителен. Легко видеть, что $R$ --- расположенное кольцо. Но $1>0$, $n\cdot1=n<x$ при любом натуральном $n$, так как $x-n>0$; значит, $R$ --- неархимедовски расположенное кольцо.
\end{rusblock}
\vocabline{\textbf{переносить} 75 \emph{carry across, transfer} (нос- fer- carry); \textbf{выполнение} 87 \emph{fulfilment} (полн- full); \textbf{обычный} 107 \emph{usual} (ук- be accustomed to).}

\readhead{B20}{Натуральные числа}{The natural numbers}
\begin{rusblock}
\textbf{Определение (Пеано).} Натуральными числами называются элементы всякого непустого множества $N$, в котором для некоторых элементов $a$ и $b$ существует отношение «$b$ следует за $a$» (число, следующее за $a$, будем обозначать через $a'$), удовлетворяющее следующим аксиомам:
\begin{enumerate}[leftmargin=2.2em]
\item существует число $1$, не следующее ни за каким числом, т.е. $a'\ne1$ для любого числа $a$;
\item для любого числа $a$ существует следующее за ним число $a'$ и притом только одно, т.е. из $a=b$ следует $a'=b'$;
\item любое число следует не более чем за одним числом, т.е. из $a'=b'$ следует $a=b$;
\item \textbf{аксиома индукции:} любое подмножество $M$ множества $N$ натуральных чисел, обладающее свойствами
  \begin{enumerate}[label=\Alph*),leftmargin=2em]
  \item $1$ принадлежит $M$;
  \item если число $a$ принадлежит $M$, то следующее число $a'$ также принадлежит $M$,
  \end{enumerate}
содержит все натуральные числа, т.е. совпадает с $N$.
\end{enumerate}
\end{rusblock}
\vocabline{\textbf{за} 52 \emph{after, in the capacity of}; \textbf{чем} 50 \emph{than}; \textbf{содержать} 105 \emph{contain} (держ- ten- hold).}

\readhead{B21}{Сложение и умножение натуральных чисел}{Addition and multiplication of natural numbers}
\begin{rusblock}
\textbf{Определение.} Сложением натуральных чисел называется такое соответствие, которое натуральным числам $a$ и $b$ сопоставляет одно и только одно натуральное число $a+b$ и обладает следующими свойствами:
\[
a+1=a',\qquad a+b'=(a+b)'.
\]
Числа $a$ и $b$ называются слагаемыми, а число $a+b$ --- их суммой.

\textbf{Теорема.} Сложение натуральных чисел существует и притом только одно. Другими словами, сложение всегда выполнимо однозначно.

\textbf{Определение.} Умножением натуральных чисел называется такое соответствие, которое натуральным числам $a$ и $b$ сопоставляет одно и только одно натуральное число $ab$ и обладает следующими свойствами:
\[
a\cdot1=a,\qquad ab'=ab+a.
\]
Числа $a$ и $b$ называются сомножителями, а число $ab$ --- произведением.
\end{rusblock}
\vocabline{\textbf{сопоставлять} 85 \emph{to associate} (ста- set).}

\readhead{B22}{Порядок натуральных чисел}{Order of the natural numbers}
\begin{rusblock}
\textbf{Определение.} Если для данных чисел $a$ и $b$ существует число $k$ такое, что $a=b+k$, то говорят, что $a$ больше $b$, и $b$ меньше $a$ ($a>b$, $b<a$).

\textbf{Теорема.} 1) Для любых чисел $a$ и $b$ имеет место одно и только одно из соотношений $a=b$, $a>b$, $b>a$. 2) Из $a>b$, $b>c$ следует $a>c$.

Другими словами, множество $N$ натуральных чисел является упорядоченным множеством.

\textbf{Теорема.} Единица --- наименьшее из натуральных чисел, т.е. $a\ge1$ для любого $a$.

\textbf{Теорема.} Если $a<b$, то $a+1\le b$.

\textbf{Теорема.} Любое непустое множество $A$ натуральных чисел содержит наименьшее число, т.е. число, меньшее всех других чисел данного множества.
\end{rusblock}

\readhead{B23}{Вычитание и деление натуральных чисел}{Subtraction and division of natural numbers}
{\noindent\bfseries\ru{Алгоритм Евклида}\par}
{\noindent\itshape Euclid's algorithm\par}
\begin{rusblock}
\textbf{Определение.} Вычитанием натуральных чисел называется операция, обратная сложению, т.е. соответствие, которое числам $a$ и $b$ сопоставляет число $a-b$ (называемое разностью $a$ и $b$) такое, что
\[
(a-b)+b=a.
\]
Делением называется операция, обратная умножению.

\textbf{Теорема (алгоритм деления с остатком).} Если $a>b$ и $a$ не делится на $b$ ($a\nmid b$), то существует одна и только одна пара чисел $q$ и $r$ такая, что
\[
a=bq+r,\qquad r<b.
\]

\textbf{Определение.} Наибольшее из чисел, делящих каждое из чисел $a,b$, называется наибольшим общим делителем (НОД) чисел $a$ и $b$. Евклиду принадлежит метод отыскания наибольшего общего делителя (НОД), дающий также и доказательство его существования.
\end{rusblock}
\vocabline{\textbf{вычитание} 104 \emph{subtraction} (чит- reckon); \textbf{остаток} 85 \emph{remainder} (ста- stand); \textbf{давать} 95 \emph{to give} (да-).}

\readhead{B24}{Фундаментальная теорема арифметики}{Fundamental theorem of arithmetic}
\begin{rusblock}
Натуральное число $p$ называется простым, если оно отлично от $1$ и не имеет делителей, кроме $p$ и $1$. Число, отличное от $1$ и не простое, называется составным. Число $1$ не считается ни простым, ни составным.

\textbf{Теорема.} Любое натуральное число $n>1$ обладает по меньшей мере одним простым делителем $p$ (наименьшее из его делителей, отличных от $1$).

\textbf{Теорема.} Для каждого простого числа $p$ существует простое число, большее $p$; таким образом, последовательность простых чисел $2,3,5,\ldots$ бесконечна.

\textbf{Доказательство.} Пусть $p_1,p_2,\ldots,p_n$ --- все простые числа $\le p$. Образуем число
\[
n'=p_1p_2\cdots p_n+1.
\]
Это число не делится ни на одно из чисел $p_1,p_2,\ldots,p_n$, но мы видели, что оно делится на некоторое простое число $q$. Значит, $q>p$.

Число $n'$ указанного вида не обязано быть простым. Например,
\[
2\cdot3\cdot5\cdot7\cdot11\cdot13+1=30031=59\cdot509.
\]

\textbf{Фундаментальная теорема арифметики.} Любое натуральное число, отличное от $1$, разлагается в произведение простых чисел и притом единственным образом до порядка сомножителей.
\end{rusblock}
\vocabline{\textbf{кроме} 52 \emph{beyond, except}; \textbf{составной} 85 \emph{composite} (ста- pos- set); \textbf{указать} 98 \emph{indicate} (каз- show); \textbf{обязать} 105 \emph{bind, compel} (вяз- knot); \textbf{разлагать} 82 \emph{expand} (лаг- lay).}

\readhead{B25}{Принцип расширения в арифметике и алгебре}{The extension principle in arithmetic and algebra}
\begin{rusblock}
Чтобы определить целые, рациональные, действительные и комплексные числа, а также многочлены и алгебраические дроби, мы последовательно построим ряд расширений множества $A$ (множества натуральных чисел), обладающих следующими свойствами по отношению к расширяемому множеству $A$:
\begin{enumerate}[leftmargin=2.2em]
\item $A$ есть подмножество множества $B$;
\item интересующие нас операции или отношения элементов множества $A$ определены также для элементов множества $B$, причём их значения для элементов $A$, рассматриваемых как элементы $B$, совпадают с теми, какие они имели до расширения в множестве $A$;
\item в множестве $B$ выполнима операция, которая в $A$ была невыполнима или не всегда выполнима;
\item расширение $B$ является минимальным из всех расширений данного $A$, обладающих свойствами 1), 2), 3), и определяется данным $A$ однозначно с точностью до изоморфизма.
\end{enumerate}
\end{rusblock}
\bcorrnote{底本首句误作“по отношению к расширяемому множеству $B$”，而紧随其后的第1项明确规定 $A\subset B$，故被扩张者应为 $A$；此处将 $B$ 改为 $A$。第2项另有“совпадают с тем, какой”这一数一致错误，按原意改为“совпадают с теми, какие”。}
\vocabline{\textbf{расширение} 90 \emph{extension} (шир- wide); \textbf{дробь} 107 \emph{fraction}; \textbf{последовательно} 103 \emph{successively}; \textbf{построить} 106 \emph{construct} (стро- stru-); \textbf{ряд} 106 \emph{row, series}; \textbf{расширять} 90 \emph{widen, extend} (шир-); \textbf{до} 52 \emph{to, up to, before}; \textbf{точность} 107 \emph{precision} (тык- pierce).}

\readhead{B26}{Выполнимость операций в расширении}{Performability of operations in an extension}
\begin{rusblock}
Объясняем на примерах требование выполнимости операции.

Для натуральных чисел не всегда выполнимо вычитание. В области целых чисел оно всегда выполнимо. Для целых чисел не всегда выполнимо деление. Для рациональных чисел оно выполнимо всегда (кроме деления на $0$).

Для рациональных чисел не всегда выполнима операция перехода к пределу фундаментальной последовательности. Для действительных чисел она всегда выполнима.

Для действительных чисел не всегда выполнима операция вычисления квадратного корня. Для комплексных чисел она всегда выполнима. Тем же принципом в алгебре расширяют одно поле, где данный многочлен не имеет корней, до поля, в котором он имеет корень, т.е. выполнима операция решения данного уравнения.
\end{rusblock}
\bcorrnote{底本这里把“переход к пределу”无条件地说成在实数中“всегда выполнима”。严格地说，这里所需的是完备性：柯西（фундаментальная）序列在实数中有极限；故仅补入“фундаментальной последовательности”这一必要条件。}
\vocabline{\textbf{выполнимость} 87 \emph{fulfillability, performability} (полн- full); \textbf{объяснять} 91 \emph{clarify} (яс- clar- clear); \textbf{требование} 106 \emph{requirement} (треб- demand); \textbf{корень} 107 \emph{root}.}

\readhead{B27}{Классы эквивалентности}{Equivalence classes}
\begin{rusblock}
\textbf{Теорема.} Если для элементов множества $M$ определено отношение эквивалентности $a\sim b$ (словами, $a$ эквивалентно $b$), т.е. отношение, обладающее следующими свойствами:
\[
a\sim a,
\]
\[
\text{если }a\sim b,\text{ то }b\sim a,
\]
\[
\text{если }a\sim b,\ b\sim c,\text{ то }a\sim c,
\]
то этим однозначно определено разбиение множества $M$ на непересекающиеся подмножества, обладающие тем свойством, что любые элементы одного и того же подмножества эквивалентны и любые элементы различных подмножеств неэквивалентны (разбиение на классы эквивалентных элементов).
\end{rusblock}

\readhead{B28}{Кольцо целых чисел до изоморфизма}{The ring of integers up to isomorphism}
\begin{rusblock}
За исходный элемент конструкции кольца целых чисел принимаем упорядоченную пару $(a,b)$ натуральных чисел. Пусть $M$ --- множество всех таких пар. Определяем отношение эквивалентности пар так:
\[
(a,b)\sim(c,d)
\]
тогда и только тогда, когда
\[
a+d=b+c.
\]
Далее, определяем сложение и умножение пар так:
\[
(a,b)+(c,d)=(a+c,b+d),
\]
\[
(a,b)(c,d)=(ac+bd,ad+bc).
\]
Пусть $C_0$ есть множество всех классов эквивалентных пар множества $M$. Суммой (произведением) двух классов $\alpha$ и $\beta$ из $C_0$ называется тот класс $\alpha+\beta$ (соответственно $\alpha\beta$), который содержит сумму (произведение) пары класса $\alpha$ и пары класса $\beta$.

\textbf{Теорема.} Множество $C_0$ с этими операциями есть кольцо.
\end{rusblock}
\vocabline{\textbf{исходный} 76 \emph{initial} (ход- go).}

\readhead{B29}{Кольцо целых чисел}{The ring of integers}
\begin{rusblock}
Но кольцо $C_0$ не удовлетворяет нашему определению кольца целых чисел, потому что не содержит натуральных чисел: его элементы --- классы эквивалентных пар натуральных чисел. Чтобы получить из $C_0$ кольцо целых чисел, необходимо включить в $C_0$ множество натуральных чисел $N$.

Любой класс $\alpha$ кольца $C_0$, отличный от нуля, состоит из пар $(a,b)$, где $a\ne b$. Пусть класс $\alpha$ называется классом первого рода, если $a>b$, т.е. если $a=b+k$, где $k$ --- натуральное число, и второго рода, если $a<b$. Пусть $N_1$ и $N_2$ --- соответственно множества классов первого и второго рода. Легко показать, что множество $N_1$ классов первого рода изоморфно множеству $N$ натуральных чисел относительно операций сложения и умножения.

Построим теперь искомое кольцо целых чисел $C$. Пусть $C$ --- множество, полученное из $C_0$ путём замены каждого класса $(a,b)$ первого рода натуральным числом $k=a-b$.

\textbf{Теорема.} Кольцо $C$, построенное выше, есть кольцо целых чисел.
\end{rusblock}
\vocabline{\textbf{наш} 45 \emph{our}; \textbf{потому что} 50 \emph{because}; \textbf{род} 106 \emph{genus, kind} (род- generation); \textbf{показать} 98 \emph{show} (каз-).}

\readhead{B30}{Поле рациональных чисел}{The field of rational numbers}
\begin{rusblock}
За исходный элемент построения поля рациональных чисел принимаем упорядоченную пару $(a,b)$ целых чисел, причём второе число отлично от нуля. Пусть $M$ --- множество всех таких пар. Определяем отношение эквивалентности пар так:
\[
(a,b)\sim(c,d)
\]
тогда и только тогда, когда
\[
ad=bc,
\]
и определяем сложение и умножение пар так:
\[
(a,b)+(c,d)=(ad+bc,bd),
\]
\[
(a,b)(c,d)=(ac,bd).
\]
Теперь за $\Gamma_0$ принимаем множество всех классов эквивалентных пар множества $M$ и за $\Gamma$ --- множество, полученное из поля $\Gamma_0$ путём замены каждого класса $(b,c)$, для которого $b$ делится на $c$, целым числом $d=b/c$.
\end{rusblock}

\readhead{B31}{Поля отношений}{Quotient fields}
\begin{rusblock}
Включение кольца целых чисел в поле рациональных чисел обобщается до включения любой области целостности в поле.

\textbf{Определение.} Полем отношений данного кольца $R$ называется минимальное множество $P$ со следующими свойствами:
\begin{enumerate}[leftmargin=2.2em]
\item $P$ содержит $R$;
\item $P$ является полем;
\item сложение и умножение в кольце $R$ совпадают с теми операциями над теми же элементами в поле $P$.
\end{enumerate}
Не все кольца обладают полями отношений; так, кольцо с делителями нуля не имеет поля отношений, так как оно не может содержаться ни в каком поле, где делители нуля отсутствуют. Таким образом, кольцо, обладающее полем отношений, является областью целостности.

\textbf{Теорема.} Для любой области целостности $R$, содержащей более одного элемента, существует поле отношений $P$ и до изоморфизма только одно.
\end{rusblock}
\vocabline{\textbf{включение} 105 \emph{inclusion} (ключ- clus- close); \textbf{какой} 45 \emph{of what kind}; \textbf{отсутствовать} 72 \emph{be absent}.}

\readhead{B32}{Поле действительных чисел}{The field of real numbers}
\begin{rusblock}
За исходный элемент построения поля действительных чисел $D$ мы принимаем фундаментальную последовательность
\[
a_1,a_2,\ldots=\{a_n\}
\]
рациональных чисел, т.е. последовательность, обладающую таким свойством: для любого рационального числа $\varepsilon>0$ существует натуральное число $n_0$ такое, что
\[
|a_p-a_q|<\varepsilon\qquad\text{при любых }p,q>n_0.
\]
Определяем отношение эквивалентности для последовательностей так:
\[
\{a_n\}\sim\{b_n\}
\]
тогда и только тогда, когда
\[
\lim_{n\to\infty}(a_n-b_n)=0,
\]
и сложение и умножение так:
\[
\{a_n\}+\{b_n\}=\{a_n+b_n\},
\]
\[
\{a_n\}\cdot\{b_n\}=\{a_nb_n\}.
\]
Пусть $D_0$ есть множество всех классов эквивалентных последовательностей.

Пусть $\Gamma$ --- подмножество классов, содержащих последовательности вида $(a,a,a,\ldots)$, где $a$ --- рациональное число.

Пусть $D$ --- множество, полученное из множества $D_0$ путём замены каждого класса $(a,a,a,\ldots)$ из $\Gamma$ рациональным числом $a$. Тогда множество $D$ есть поле действительных чисел.
\end{rusblock}

\readhead{B33}{Поле комплексных чисел}{The field of complex numbers}
\begin{rusblock}
Пусть $K_0$ есть множество всех упорядоченных пар вида $(a,b)$, где $a$ и $b$ --- действительные числа. Сложение и умножение в множестве $K_0$ определяем по формулам:
\[
(a,b)+(c,d)=(a+c,b+d),
\]
\[
(a,b)(c,d)=(ac-bd,ad+bc).
\]
Пусть $D'$ --- множество всех пар множества $K_0$ вида $(a,0)$ и $D$ --- множество, полученное из $D'$ путём замены каждой пары $(a,0)$ из $D'$ действительным числом $a$.
\end{rusblock}

\endgroup
